El març del 2021 es va llançar el primer nanosatèl·lit de la Generalitat de Catalunya a l'espai. Els nanosatèl·lits tenen com a funció millorar les comunicacions, controlar els cabals dels cursos d'aigua i prevenir incendis. També contribueixen a la recerca i realització de missions espacials més àgils i econòmiques. Aquests nanosatèl·lits acostumen a orbitar a uns 500~km d'altura (distància respecte a la superfície de la Terra).

\begin{apartat}{1,25}
Suposant que l'òrbita d'un d'aquests nanosatèl·lits és circular, a partir de la llei de la gravitació universal deduïu-ne l'expressió de la velocitat orbital en funció del radi orbital. Calculeu també la velocitat i el període orbitals d'aquests nanosatèl·lits.
\end{apartat}

\begin{apartat}{1,25}
Partint de la llei de la conservació de l'energia mecànica (negligiu la força de fregament), deduïu l'expressió de la velocitat de llançament necessària per a posar en òrbita un satèl·lit en funció del radi orbital. Calculeu la velocitat de llançament necessària per a posar en òrbita el nanosatèl·lit a 500~km d'altura.
\end{apartat}

\dades{%
  $G = 6,67\times 10^{-11}\ \mathrm{N\,m^2\,kg^{-2}}$.\\
  $M_\mathrm{Terra} = 5,98\times 10^{24}\ \mathrm{kg}$.\\
  $R_\mathrm{Terra} = 6,37\times 10^{3}\ \mathrm{km}$.}
